\subsection{矩阵系数与 G-有限函数}

命题 6. --- \(\mathrm{L}^{2}(\mathrm{G})\) 的下列子空间相等：

\begin{enumerate}
\def\labelenumi{\alph{enumi})}
\item
  空间 \(\Theta (\mathbf{G})\)
\item
  \(\varrho_{\mathrm{G}}(\pi)\) 的 \(\mathrm{L}^2 (\mathrm{G})\) 子空间的代数直和。
\item
  \(\gamma_{\mathrm{G}}\) 的 G-有限向量空间（参见 V, p. 376）；
\item
  \(\delta_{G}\) 的 G-有限向量空间；
\item
  \((\mathbf{G} \times \mathbf{G})\)-有限向量空间的 \(\varrho_{\mathrm{G}}\)。
\end{enumerate}

特别地，\(\gamma_{\mathrm{G}}\)、\(\delta_{G}\) 或 \(\varrho_{\mathrm{G}}\) 的每个 G-有限向量都属于 \(\mathcal{C}(\mathrm{G})\)。

以 \(F_{a}\)（分别以 \(F_{b}\)、\(F_{c}\)、\(F_{d}\)、\(F_{e}\)）表示由条件 a)（分别由 b)、c)、d)、e)）定义的空间。我们已经注意到 \(F_{a} = F_{b}\)。

有 \(F_{b} \subset F_{c}\)，因为 \(\varrho_{\mathrm{G}}(\pi)\) 是 \(\gamma_{G}\) 的有限维子表示，对所有 \(\pi \in \widehat{G}\) 成立。反之，设 f 是 \(\gamma_{G}\) 的 G-有限向量。子空间 \(E_{f}\) 由函数 \(\gamma_{\mathrm{G}}(g)f\)（其中 \(g \in G\)）生成，是 \(\gamma_{G}\) 的有限维子表示。它等于其 \(\pi\)-同型分量的直和，对所有 \(\pi \in \widehat{G}\)（V, p. 457, 命题 1）。根据 V, p. 463, 推论 4，这意味着 \(E_{f}\) 包含于各空间 \(\varrho_{\mathrm{G}}(\pi)\) 的和中，从而 \(F_{c} \subset F_{b}\)。

用类似论证可得 \(F_{b} = F_{d}\)；由于 \(\varrho_{\mathrm{G}}(\pi)\) 是 \(\varrho_{G}\) 的子表示，同样可证明 \(F_{b} = F_{e}\)。

推论。--- 设 \(E \subset L^{2}(G)\) 是定义 \(\gamma_{G}\)（分别为 \(\delta_{G}\)、\(\varrho_{G}\)）的一个子表示的有限维向量子空间。则 E 包含于 \(\mathcal{C}(G)\)。

事实上，E 的每个元都是表示 \(\gamma_{G}\) 的 G-有限向量。

